\section{Wyniki}

% \subsection{Pierwsza próba}
W pierwszych próbach, czyli przy pierwszym kontakcie pracownika z danym pytaniem, zarejestrowano
\FirstN{} odpowiedzi \NEmployees{} pracowników. Poprawność wynosiła \FirstAcc\% (95\% PU
\FirstAccCI; tab.~\ref{tab:pierwsza_proba_wiek}). W przedziałach wieku mieściła się w zakresie od
\FirstAccMin\% (\FirstAccMinGroup{} lat) do \FirstAccMax\% (\FirstAccMaxGroup{} lat), a przedziały
ufności w grupach o mniejszej liczbie pracowników były szerokie. Po uwzględnieniu trudności pytań
szanse poprawnej odpowiedzi w pierwszej próbie rosły wraz z wiekiem (OR \FirstAgeOR{} na 10 lat,
95\% PU \FirstAgeORCI; $p \FirstAgeP$). Zależność ta nie miała jednak charakteru monotonicznego,
ponieważ najniższą poprawność odnotowano w przedziale \FirstAccMinGroup{} lat, a nie w przedziale
najmłodszym. % SPRAWDŹ: kierunek OR i położenie minimum po zmianie danych

\begin{table}[H]
\centering
\caption{Poprawność w pierwszej próbie ogółem i według przedziału wieku pracownika.}
\label{tab:pierwsza_proba_wiek}
\footnotesize
\setlength{\tabcolsep}{4pt}
\tabelacsv{l*{7}{r}}{pierwsza_proba_wiek}% 7 = ogółem + przedziały wieku + brak danych
\uwagatabeli{Pierwsza próba to pierwsza odpowiedź w sekwencji. 95\% PU wyznaczono metodą bootstrap.
Zależność od wieku oszacowano regresją logistyczną poprawności pierwszej próby względem środka
przedziału wieku (na 10 lat) z efektami stałymi pytania i błędami standardowymi klastrowanymi względem
pracowników (\FirstAgeN{} odpowiedzi \FirstAgeEmp{} pracowników o znanym wieku): OR \FirstAgeOR,
95\% PU \FirstAgeORCI, $p \FirstAgeP$.}
\end{table}

\subsection{Zmiany w sekwencji (PB1)}
Poprawność rosła wraz z numerem próby: z \AccAttOne\% \AccAttOneCI{} w pierwszej próbie do
\AccAttMax\% \AccAttMaxCI{} w próbie \AttMax{}. i kolejnych (rys.~\ref{fig:krzywa_uczenia}).
Liczba sekwencji malała jednak wraz z numerem próby, dlatego zmianę
oceniono również w kohorcie \CohortN{} sekwencji. W kohorcie poprawność wzrosła z \CohortAccOne\%
\CohortAccOneCI{} w próbie~1 do \CohortAccK\% \CohortAccKCI{} w próbie~\CohortK{}, czyli
o \SeqChange{} p.p. (95\% PU \SeqChangeCI; tab.~\ref{tab:sekwencje_wiek}). Wzrost obserwowany w tych
samych sekwencjach wskazuje, że poprawa nie wynika z wcześniejszego kończenia sekwencji z niższą
poprawnością. Odsetek błędów był funkcją potęgową numeru próby (wykładnik \PowerLawB,
$R^2 = \PowerLawRsq$), zgodnie z potęgowym prawem praktyki \citep{newell1981}.

Wzrost poprawności w kohorcie wystąpił we wszystkich przedziałach wieku i wynosił od \SeqChangeMin{}
p.p. (\SeqChangeMinGroup{} lat) do \SeqChangeMax{} p.p. (\SeqChangeMaxGroup{} lat;
tab.~\ref{tab:sekwencje_wiek}). Tempo poprawy nie zależało od wieku: iloraz szans dla interakcji
numeru próby z wiekiem wynosił \SeqAgeOR{} na 10 lat (95\% PU \SeqAgeORCI; $p \SeqAgeP$). Różnice
między przedziałami wieku dotyczyły zatem raczej poziomu wyjściowego niż przyrostu poprawności
w kolejnych próbach. % SPRAWDŹ: znak zmiany we wszystkich przedziałach i istotność interakcji

\begin{figure}[H]
\centering
\includegraphics{\FigDir rys_krzywa_uczenia}
\caption{Poprawność według numeru próby we wszystkich sekwencjach dostępnych w danej próbie (ich liczba
maleje wraz z numerem próby). Słupki błędów: 95\% PU (bootstrap). Linia przerywana: dopasowanie funkcji
potęgowej do odsetka błędów. Próby \MaxAttemptPooled{}. i dalsze połączono. Przebieg w kohorcie
przedstawiono na rys.~\ref{fig:stany_kohorty}.}
\label{fig:krzywa_uczenia}
\end{figure}

\begin{table}[H]
\centering
\caption{Poprawność w kolejnych próbach w kohorcie, ogółem i według przedziału wieku pracownika.}
\label{tab:sekwencje_wiek}
\footnotesize
\setlength{\tabcolsep}{4pt}
\tabelacsv{l*{7}{r}}{sekwencje_wiek}% 7 = ogółem + przedziały wieku + brak danych
\uwagatabeli{Kohorta: sekwencje liczące co najmniej \CohortK{} prób, próby od 1 do \CohortK{}.
p.p.: punkty procentowe; 95\% PU zmiany wyznaczono metodą bootstrap. Zależność tempa poprawy od wieku
oszacowano regresją logistyczną poprawności względem $\log_2(\text{numer próby})$ i jego interakcji
z wiekiem (środek przedziału, na 10 lat), z efektami stałymi pracownika i pytania oraz błędami
standardowymi klastrowanymi względem pracowników (\SeqAgeN{} odpowiedzi pracowników o znanym wieku).}
\end{table}

\subsection{Wynik poprzedniej próby, odstęp i zapominanie (PB2)}
Wzrost poprawności opisany w PB1 jest wynikiem netto dwóch procesów zachodzących między kolejnymi
próbami: korygowania błędów i zapominania odpowiedzi poprawnych. Aby je rozdzielić, każdą próbę od
drugiej wzwyż powiązano z wynikiem próby poprzedniej i odstępem od niej. Harmonogram powtórek nie
zniekształcał tego porównania, ponieważ pytanie wracało po odpowiedzi poprawnej i błędnej z podobnym
prawdopodobieństwem (\ReaskAfterCorrect\% i \ReaskAfterIncorrect\%) i po podobnym odstępie (mediana
\GapAfterCorrect{} i \GapAfterIncorrect{} dni). Odpowiedź poprawna powtórzyła się w kolejnej próbie
w \RetainPct\% \RetainCI{} przypadków, a odpowiedź błędna została skorygowana w \RecoverPct\%
\RecoverCI{} (tab.~\ref{tab:przejscia}). Korekty błędów przeważały zatem nad zapominaniem
(\ForgetPct\%), co tłumaczy przyrost poprawności w kolejnych próbach. Upływ czasu działał
w przeciwnym kierunku, lecz tylko w odniesieniu do odpowiedzi poprawnych: odsetek ich powtórzeń
malał z \RetShortCorrect\% przy odstępie krótszym niż 14 dni do \RetLongCorrect\% przy odstępie
wynoszącym co najmniej 240 dni, podczas gdy odsetek skorygowanych błędów nie wykazywał wyraźnego
trendu (\RetShortIncorrect\% i \RetLongIncorrect\% w tych samych przedziałach; rys.~\ref{fig:odstep}).
Wiedza potwierdzona poprawną odpowiedzią podlegała więc
zapominaniu między powtórkami \citep{cepeda2006,cepeda2008}, natomiast szansa skorygowania błędu po
informacji zwrotnej nie zależała od czasu, który upłynął do kolejnego kontaktu z pytaniem. % SPRAWDŹ: brak trendu dla odpowiedzi błędnych po zmianie danych

\begin{figure}[H]
\centering
\includegraphics{\FigDir rys_odstep}
\caption{Poprawność według odstępu, osobno dla prób poprzedzonych odpowiedzią poprawną (koła)
i błędną (kwadraty). Słupki błędów: 95\% PU (bootstrap).}
\label{fig:odstep}
\end{figure}

\begin{table}[H]
\centering
\caption{Przejścia między kolejnymi próbami tego samego pytania.}
\label{tab:przejscia}
\small
\tabelacsv{lrrlr}{przejscia}
\uwagatabeli{Próby od drugiej wzwyż według wyniku próby poprzedniej. 95\% PU wyznaczono metodą
bootstrap.}
\end{table}

\subsection{Utrwalone błędy (PB3)}
Większość błędów miała charakter przejściowy: w \EverCorrectedPct\% z \FirstWrongSeqN{} sekwencji
rozpoczętych odpowiedzią błędną pracownik odpowiedział poprawnie w którejś z kolejnych prób. Szansa
korekty malała jednak z każdym kolejnym błędem. Po jednym błędzie poprawność wynosiła \CorrRunOne\%
\CorrRunOneCI, po dwóch błędach z rzędu \CorrRunTwo\%, po trzech \CorrRunThree\%, a po co najmniej
\RunMax{} błędach \CorrRunMax\% \CorrRunMaxCI{} (tab.~\ref{tab:seria_bledow}). Każde kolejne
wyświetlenie poprawnej odpowiedzi przynosiło więc mniejszy efekt. Spadek ten może częściowo wynikać
z tego, że wiele błędów z rzędu występowało w trudniejszych pytaniach, nie musi zatem w całości odzwierciedlać
utrwalania się błędu.

O naturze powtórzonych błędów świadczy wybór dystraktora. Gdy pracownik pomylił się w dwóch kolejnych
próbach, w \SameDistractorPct\% \SameDistractorCI{} przypadków ($n = \SameDistractorN$) wskazał ten sam
dystraktor, wobec \SameDistractorExpected\% oczekiwanych przy wyborze losowym
i \SameDistractorExpectedPop\% oczekiwanych przy uwzględnieniu popularności poszczególnych dystraktorów
w danym pytaniu. Powtórzony błąd był zatem wyborem konkretnej odpowiedzi, a nie zgadywaniem. Tendencja
ta nie słabła z upływem czasu: do tego samego dystraktora wracano w \SameDistrGapShort\% przypadków przy
odstępie krótszym niż 30 dni, w \SameDistrGapMid\% przy odstępie od 30 do 89 dni i w \SameDistrGapLong\%
przy odstępie wynoszącym co najmniej 90 dni. Powtórzone błędy skupiały się przy tym w niewielkiej części
materiału, ponieważ na \PersistTopItems{} spośród \PersistItems{} pytań, w których wystąpiły, przypadało
\PersistTopItemsShare\% wszystkich powtórzonych błędów.

Zestawienie wyników PB2 i PB3 ujawnia asymetrię. Odpowiedzi poprawne słabły wraz z odstępem, natomiast
błędne przekonanie, które przetrwało pierwszą korektę, utrzymywało się niezależnie od upływu czasu
i w coraz mniejszym stopniu ulegało informacji zwrotnej. Wyświetlenie poprawnej odpowiedzi wystarczało
do skorygowania większości błędów, ale nie usuwało przyczyny, dla której pracownik wybierał określony
dystraktor. % SPRAWDŹ: monotoniczny spadek korekty i brak trendu dystraktora wg odstępu po zmianie danych

\begin{table}[H]
\centering
\caption{Poprawność według liczby błędów popełnionych z rzędu bezpośrednio przed daną próbą.}
\label{tab:seria_bledow}
\footnotesize
\setlength{\tabcolsep}{5pt}
\tabelacsv{l*{5}{r}}{seria_bledow}% 5 = 0, 1, 2, 3, >=RunMax
\uwagatabeli{Próby od drugiej wzwyż. Kolumna 0: próby poprzedzone odpowiedzią poprawną. 95\% PU
wyznaczono metodą bootstrap.}
\end{table}

\subsection{Podsumowanie wyników}
Wyniki pozwalają odpowiedzieć na trzy pytania badawcze.
\begin{enumerate}[label=(PB\arabic*)]
  \item Poprawność rosła wraz z numerem próby, również w kohorcie, w której zwiększyła się
  o \SeqChange{} p.p. między próbą~1 a próbą~\CohortK{}. Przyrost nie zależał od wieku pracownika,
  a wiek wiązał się jedynie z poprawnością w pierwszej próbie.
  \item Przyrost poprawności wynikał z przewagi korekt błędów (\RecoverPct\%) nad zapominaniem
  (\ForgetPct\%). Zapominanie nasilało się wraz z odstępem, natomiast odsetek korekt błędów od odstępu
  nie zależał.
  \item Większość błędów korygowano (\EverCorrectedPct\% sekwencji rozpoczętych odpowiedzią błędną),
  lecz powtórzone błędy polegały na wyborze tego samego dystraktora (\SameDistractorPct\%), nie słabły
  wraz z odstępem, a szansa ich korekty malała z każdym kolejnym błędem z rzędu, do \CorrRunMax\%
  po co najmniej \RunMax{} błędach.
\end{enumerate}
Łącznie wyniki wskazują, że program utrwalał wiedzę sprawdzaną przez pytania, przy czym wiedza ta
wymagała podtrzymywania kolejnymi powtórkami, a niewielka część pytań generowała błędy odporne na
informację zwrotną ograniczoną do wyświetlenia poprawnej odpowiedzi. % SPRAWDŹ
Rys.~\ref{fig:stany_kohorty} zestawia te wyniki w kohorcie: linia poprawności rośnie (PB1), korekty
błędów przeważają nad zapominaniem (PB2), a powtórzone błędy utrzymują się we wszystkich kolejnych
próbach (PB3).

\begin{figure}[H]
\centering
\includegraphics{\FigDir rys_stany_kohorty}
\caption{Poprawność i stan odpowiedzi w kohorcie (\CohortN{} sekwencji, próby od 1 do \CohortK{}).
Słupki: udział odpowiedzi według wyniku bieżącej i poprzedniej próby (definicje IX--XI
w podrozdziale~\ref{sec:definicje}); w próbie~1 tylko odpowiedzi poprawne i błędne. Linia: poprawność z 95\% PU (bootstrap),
wartości nad słupkami.}
\label{fig:stany_kohorty}
\end{figure}

% \section{Wyniki - stare do usuniecia}

% \subsection{Próba i harmonogram powtórek}
% Zbiór danych obejmował \NResponses{} odpowiedzi \NEmployees{} pracowników na \NItems{} pytań
% w \NModules{} modułach, zebranych od \DateStart{} do \DateEnd{} (tab.~\ref{tab:zbior_danych}).
% Mediana liczby odpowiedzi na pracownika wynosiła \RespPerEmpMedian{} (IQR
% \RespPerEmpQOne--\RespPerEmpQThree), udzielonych na \ItemsPerEmpMedian{} różnych pytań w ciągu
% \ActiveDaysMedian{} dni aktywności, a mediana czasu uczestnictwa w programie wynosiła
% \SpanMonthsMedian{} miesięcy. Okresy bez odpowiedzi odpowiadają zawieszeniom programu na czas
% przeglądu pytań (rys.~\ref{fig:aktywnosc}). Wśród \DemoMatched{} pracowników o znanych danych
% demograficznych kobiety stanowiły \DemoWomenPct\%, a mediana wieku wynosiła \AgeMedian{} lat
% (tab.~\ref{tab:demografia}). Ogólny odsetek poprawnych odpowiedzi wynosił \AccOverall\%, a w pierwszych
% próbach \AccFirst\%.

% Harmonogram powtórek nie zależał od wyniku poprzedniej próby: prawdopodobieństwo kolejnej próby
% wynosiło \ReaskAfterCorrect\% po odpowiedzi poprawnej i \ReaskAfterIncorrect\% po błędnej, natomiast
% mediana odstępu do kolejnej próby wynosiła odpowiednio \GapAfterCorrect{} i \GapAfterIncorrect{} dni.
% Kontynuacja sekwencji nie była więc selekcjonowana według wyników. % SPRAWDŹ

% \subsection{Uczenie się w kolejnych próbach (PB1)}
% Odsetek poprawnych odpowiedzi rósł wraz z numerem próby: z \AccAttOne\% \AccAttOneCI{} w pierwszej
% próbie do \AccAttMax\% \AccAttMaxCI{} w próbie \AttMax{}. i kolejnych (rys.~\ref{fig:krzywa_uczenia},
% tab.~\ref{tab:krzywa_uczenia}). Wzrost był podobny w zbalansowanej kohorcie \CohortN{} sekwencji
% liczących co najmniej \CohortK{} prób (\CohortAccOne\% \CohortAccOneCI{} w próbie~1 wobec \CohortAccK\%
% \CohortAccKCI{} w próbie~\CohortK), nie wynika więc z selektywnej kontynuacji przez pracowników
% osiągających lepsze wyniki. Odsetek błędów był funkcją potęgową numeru próby (wykładnik \PowerLawB,
% $R^2 = \PowerLawRsq$), zgodnie z potęgowym prawem praktyki \citep{newell1981}.

% W modelu uczenia M1 każde podwojenie numeru próby wiązało się z \ORMOneAttempt-krotnie wyższymi
% szansami poprawnej odpowiedzi (95\% PU \ORMOneAttemptCI; tab.~\ref{tab:regresja}). Efekt ten wyjaśniał
% jednak niewielką część zmienności poprawności (brzegowy $R^2 = \RsqMarginal$ wobec warunkowego
% $R^2 = \RsqConditional$), a pytania odpowiadały za większą jej część niż pracownicy (ICC
% \ICCItem{} wobec \ICCEmployee). Na poziomie pracowników wynik w trzeciej próbie był lepszy niż
% w pierwszej u \EmpImproved{} z \EmpN{} osób, a gorszy u \EmpWorsened{} (średnia zmiana
% \EmpMeanChange{} p.p., 95\% PU \EmpMeanChangeCI; tab.~\ref{tab:zmiana_pracownikow}).

% \begin{table*}[htbp]
% \centering
% \caption{Modele regresji logistycznej poprawności odpowiedzi: ilorazy szans [95\% PU].}
% \label{tab:regresja}
% \small
% \setlength{\tabcolsep}{3.5pt}
% \tabelacsv{llll}{regresja}
% \uwagatabeli{Definicje modeli M1--M3 podano w podrozdziale ,,Modele statystyczne''. Wszystkie modele
% zawierają efekty stałe pracownika i pytania; błędy standardowe klastrowane względem pracowników.
% Pracowników i pytania bez zmienności poprawności usunięto przed estymacją (M1: \MOneDropped{}
% odpowiedzi). OR dla numeru próby i odstępu dotyczą podwojenia wartości, dla stażu jednego roku.}
% \end{table*}

% \subsection{Wynik poprzedniej próby, odstęp i zapominanie (PB2)}
% W modelu retencji M2, po uwzględnieniu numeru próby, poprawna poprzednia odpowiedź wiązała się
% z wyższymi szansami poprawnej odpowiedzi (OR \ORMTwoPrev, 95\% PU \ORMTwoPrevCI), natomiast dłuższy
% odstęp od poprzedniej próby wiązał się z niższymi szansami (OR \ORMTwoGap{} na podwojenie odstępu,
% 95\% PU \ORMTwoGapCI; tab.~\ref{tab:regresja}). Wśród odpowiedzi poprzednio poprawnych odsetek
% poprawnych zmniejszał się z \RetShortCorrect\%, gdy pytanie powtórzono w ciągu 14 dni, do
% \RetLongCorrect\% po 240 dniach lub dłużej; odpowiedzi poprzednio błędne były korygowane odpowiednio
% w \RetShortIncorrect\% i \RetLongIncorrect\% przypadków (rys.~\ref{fig:odstep},
% tab.~\ref{tab:retencja_wg_odstepu}). Wzorzec ten jest zgodny z zapominaniem między powtórkami
% \citep{cepeda2006,cepeda2008}. % SPRAWDŹ

% \subsection{Przejścia i utrwalone błędy (PB3)}
% Spośród odpowiedzi poprzednio poprawnych \RetainPct\% \RetainCI{} było poprawnych również w kolejnej
% próbie, a \ForgetPct\% zapomniano. Spośród odpowiedzi poprzednio błędnych \RecoverPct\% \RecoverCI{}
% skorygowano, a \PersistPct\% pozostało błędnych, mimo że po każdej odpowiedzi wyświetlano poprawną
% odpowiedź (tab.~\ref{tab:przejscia}). Gdy błąd się powtarzał (\SameDistractorN{} przypadków),
% pracownicy wybierali ten sam dystraktor w \SameDistractorPct\% \SameDistractorCI{} przypadków, wobec
% \SameDistractorExpected\% oczekiwanych przy wyborze losowym i \SameDistractorExpectedPop\% oczekiwanych
% przy uwzględnieniu popularności poszczególnych dystraktorów. Wskazuje to na utrwalone błędne
% przekonania, a nie na zgadywanie. % SPRAWDŹ

% \subsection{Odporność wyników i nowe pytania}
% Oszacowanie efektu numeru próby pozostawało zbliżone we wszystkich specyfikacjach (OR od \ORRange;
% rys.~\ref{fig:wrazliwosc}, tab.~\ref{tab:wrazliwosc}) i wynosiło \ORSensCohort{} w zbalansowanej
% kohorcie, \ORSensLenThree, \ORSensLenFive{} i \ORSensLenEight{} w sekwencjach liczących 3--4, 5--7
% i co najmniej 8 prób, \ORSensGee{} w uśrednionym populacyjnie modelu GEE oraz \ORSensMixed{} w modelu
% mieszanym. Oszacowanie było wyższe po dodaniu stażu jako predyktora (\ORSensTenure), ponieważ staż
% jest skorelowany z numerem próby.

% W modelu nowych pytań M3 odsetek poprawnych odpowiedzi przy pierwszym kontakcie z pytaniem nie rósł
% wraz ze stażem (OR \ORMThreeTenure{} na rok, 95\% PU \ORMThreeTenureCI; tab.~\ref{tab:regresja}).
% Przedział ufności obejmuje zarówno brak efektu, jak i niewielki spadek, dlatego wynik oznacza brak
% dowodu na uogólnienie poprawy na nowe pytania, a nie dowód jego braku. % SPRAWDŹ

% \begin{figure}[tbp]
% \centering
% \includegraphics{\FigDir rys_wrazliwosc}
% \caption{Iloraz szans (95\% PU) poprawnej odpowiedzi na podwojenie numeru próby w różnych
% specyfikacjach modelu M1. Romb: model główny; linia pionowa: brak efektu.}
% \label{fig:wrazliwosc}
% \end{figure}
